\chapter[作家的职责不只是编排（怎么可能会后悔）]{作家的职责不只是编排\protect\footnotemark（怎么可能会后悔）}

\footnotetext{译注：``The role of the writer is not simply to arrange Being according to his own lights; he must also serve as a medium to Being and remain open to its often unfathomable dictates. This is the only way the work can transcend its creator and radiate its meaning further than the author himself can see or perceive.''

“作家的职责并非只是按照自己的眼光安排存在；他还必须成为存在的媒介，向存在那些往往深不可测的律令保持敞开。唯有如此，作品才能超越其创作者，将意义辐射到作者本人所能看见或感知的范围之外。”
\par
这一点在《魔法少女小圆》中得到了充分体现。例如，小圆从最初的怯懦与依赖，逐渐接受自我牺牲的命运，最终选择挑战宿命，这一转变展现了她内心的复杂与成长；杏子从一开始的冷酷无情，到最终为他人牺牲，也反映出她在困境中作出的并非理性算计，而是由情感推动的选择。少女们的选择和行动并非简单反应，也并不遵循固定逻辑，而是充满冲突、矛盾和不可预见性，保留着某种“深不可测的律令”。
\par
丘比作为外部力量，固然设定了魔法少女的命运框架；然而这些角色在面对困境时，并不按预定轨迹行走。她们的选择、牺牲与反抗展现了她们作为独立个体的复杂性，也推动故事向不可知的方向发展，使作品获得更高层次的意义与冲突。
\par
可以说，《魔法少女小圆》不仅是一出悲剧，也是一部充满哲学反思的深刻作品。它探讨命运、牺牲、自我认知与自由意志，呈现角色如何凭借那些“深不可测”的行动，使作品本身的意义超越原初构想，变得更为深邃有力。}

第五集开头，第二篇章延续了对第一篇章的镜像。正如第一集和第四集都承担引介功能，第二集和第五集同样是在安置角色，并探查这个新世界的性质。说得客气些，这是一集几乎无事发生的剧集。

本集以一段回忆开场：丘比与沙耶香举行仪式，使她变成魔法少女；时间大概就在第四集中沙耶香向恭介保证魔法确实存在之后。尽管自系列开始以来，丘比始终显得阴森可怖，但这段回忆是截至目前最长的一次连续描写，将它呈现为一个（在这里几乎是字面意义上的）阴影人物，而构图与光线都强烈指向一场死亡戏。正如焰后来所说，沙耶香的命运在此已经锁定；实际上，她已经是个行走的死人。（几集之后我们会知道，这句话同样相当字面。）此外，镜头把丘比置于浓重阴影之中，又让一株高大的植物立在它身后，使若干画面看起来仿佛丘比长着多条尾巴，在视觉上近似强大的九尾灵狐\footnote{译注：中国神话生物，《山海经·南山经》云：“青丘之山，有兽焉，其状如狐而九尾。”至少在宋代以前，九尾狐多属祥兽；约宋代以后，九尾狐才开始偏向魅惑、欺骗等形象。及至明代小说《封神演义》成书后，九尾狐作为恶兽的观念已大致定型，狐狸亦成为狡猾的代名词。在《封神演义》中，九尾狐奉女娲之命前去诱惑商纣王，化身妲己，做了许多坏事，使纣王失去民心与江山，最终引来周武王伐纣。}，一种（通常带有恶意的）诡计精怪\cite{ref32}。

我们也第一次看到灵魂宝石如何形成。仪式强烈暗示，丘比实际上是从少女们的心脏中把灵魂宝石拉出来，以她们体内原本已经存在的某种东西塑造成型。这与前几集中麻美和丘比所说的“虽然小圆尚未成为魔法少女，却能感觉到她体内蕴藏巨大力量”相互印证，也解释了丘比为什么不直接凭自己的力量达成目标：它的力量尽管巨大，却极其受限。它可以实现愿望，充当心灵感应的中继，控制谁能看见它，并且（正如我们稍后将在系列中了解到的）同时存在于多个地点，却无法像魔法少女和魔女那样真正施展魔法、改变现实。它的几句台词甚至暗示，它所能实现的愿望取决于许愿的魔法少女本身拥有多强的力量。考虑到后几集中有人提到沙耶香并不是很强大的魔法少女，她之所以只许愿治好恭介的手，而不是治愈他的整个身体，可能正是因为她没有足够的力量治愈其余部分。

换言之，丘比是一个促成者。它使准魔法少女能够触及她们体内已经存在的力量，好让她们替它猎杀魔女，或替它彼此厮杀。正如我们在本集中所见，它全然乐于制造冲突：明知杏子即将到来，仍赋予沙耶香力量；一边向杏子提供情报，一边让沙耶香蒙在鼓里，因为她们之间的战斗正符合它的目的。

由于丘比的操控，杏子接过了焰在第一篇章中的位置，成为道德立场可疑的敌对魔法少女。杏子正是麻美曾经警告过的那类人：非常愿意与其他魔法少女战斗，只关心击败魔女所得的报酬，毫不在意保护见泷原的民众。她愿意放任使魔杀人，直到使魔成长为魔女；再加上她关于食物链的评论，以及她自己不停进食的习惯，这些都共同暗示，杏子把“吃”视为力量的表达，并围绕这种力量信奉“强权即公理”的哲学。与她相对，沙耶香则接过了麻美的角色，成为“善”的魔法少女：她为保护他人而战，并相信强者负有守护弱者的义务。

焰出手打断二人的战斗，也揭示出她已经接过麻美在第一篇章中扮演的另一重角色。焰不再试图抹去传统魔法少女结构，并以《魔法少女小圆》终将成为的那部作品取而代之；这件事已经完成。相反，她如今试图阻止故事按逻辑推进的下一步，也就是片名所预示的“魔法少女小圆”的出现。当小圆恳求她帮助沙耶香时，焰拒绝了；但当唯一的替代方案是让小圆成为魔法少女时，焰便别无选择，只能介入。

于是只剩下小圆与丘比。小圆做出了一个有趣而有意为之的选择：她不改变自己的角色。正如她曾是麻美身边没有魔力的帮手与知己，她也愿意对沙耶香如此。即便她理所当然地害怕成为魔法少女，也仍愿意冒生命危险站在沙耶香身边；而正如本集结尾所示，若有必要，她甚至愿意自己成为魔法少女。丘比同样没有改变角色，而是越来越清楚地向观众显露它的角色究竟是什么：它从一个道德和归属都成疑的角色，逐渐变为明显操控他人的人物，并逐渐被置于敌对一方。它主动协助杏子，又向小圆和沙耶香隐瞒杏子的存在。

但这个角色究竟是什么？本集早些时候，沙耶香带恭介上到屋顶，他的家人把小提琴还给他，而他演奏了《圣母颂\footnote{译注：即古诺《圣母颂》，由法国作曲家夏尔-弗朗索瓦·古诺创作，伴奏取自 J.S. 巴赫的《十二平均律钢琴曲集》。}》。这一幕之所以值得注意，是因为场景中的诸多元素不断将恭介与丘比对照起来。首先，这是本集开头沙耶香立约成为魔法少女的同一地点；尽管此时二人位于同心花环之外，而非花环之内，沙耶香和恭介的相对位置仍与回忆中沙耶香和丘比的位置相同。其次，在某些镜头里，恭介以剪影形式出现，正如回忆中的丘比；而在沙耶香视角的镜头中，恭介挡住的正是回忆里丘比所站的位置。最微妙、也最有助于理解恭介-丘比对位关系功能的一点是：丘比伸手探入沙耶香的心脏，使她成为魔法少女；而恭介的音乐也抵达了沙耶香的心，使她萌生对他的感情，并因此选择成为魔法少女。当然，接下来几集中我们会看到，恭介对沙耶香颇为迟钝且缺乏体贴，因此他与丘比还有一个共同点：二者都促成了沙耶香成为魔法少女，却都很少在意她的感受。

临近这场戏的结尾，恭介放下小提琴，而《圣母颂》仍作为背景音乐延续至整场戏结束，从叙事内声音（也就是剧中世界里的角色能听到的声音）转为外叙事声音。到目前为止，剧中一直把这种跨越叙事边界的能力呈现为魔女才拥有的能力。为了传达魔女的异界性，魔女结界通常拥有独立的艺术风格，与剧集常规画风截然不同；因此，我们会看到角色注意到这种画风变化，并对此产生恐惧反应——这是叙事内人物对外叙事效果作出的回应。此刻，恭介则从相反方向越过这道界线，成为仅有的三位跨过这一门槛的非魔女角色之一，而且是唯一一个没有任何明显“魔法”能力却能做到这一点的角色。

这是因为他的音乐是情感的表达，因而也是魔法；正如我们将在系列后期看到的，人类情感是本剧所有魔法的源泉。由此可知，人类艺术本质上也具有魔法性，艺术表达能够重塑现实。这里确实如此：剧情强烈暗示，恭介的音乐是沙耶香对他产生爱慕的源头，而这正是她成为魔法少女的原因。因此，魔法少女沙耶香每一次对物理法则的违背，都可以追溯为恭介音乐的结果。

可是，如果恭介能够在叙事内与外叙事之间跨界，而这一幕又让他与丘比形成强烈对照，那么是否意味着丘比也能做到同样的事？答案是肯定的：它正是三位跨过这一边界的非魔女角色中的第二位。丘比是一个栖居于故事内部的外叙事实体。

不妨想想：丘比创造魔法少女来服务自己的目的，明知她们会受苦，甚至正是依赖这种痛苦。它想让小圆成为魔法少女，于是把自己的一切行动都导向这一最终结果，因为它有一个小圆可以帮它解决的问题。它安排杏子与沙耶香战斗，同样出于这个理由；除非影响本身有助于它的目标，否则它丝毫不在意这会如何影响她们。那么，它的目标是什么？让它所创造的魔法少女经历悲剧与绝望。

在这一点上，它很容易令人联想到另一部魔法少女作品《萩萩公主》中的反派多罗斯玛亚\cite{ref33}。在那部作品中，多罗斯玛亚最终被揭示为一名作家，拥有将自己的创作赋予生命的能力；他利用这种能力，将一座小镇困在悲剧的循环之中\footnote{《萩萩公主》第21话：编织者们$\sim$无言歌$\sim$}。然而，那只被他变成题名角色的鸭子以纯粹的牺牲终结了悲剧循环，迫使故事走向更幸福的结局，也将多罗斯玛亚逐出这个世界，送往其他“世界/文字”。将这一情节与我们上一章讨论过的虚淵\, 玄后记相对照：“我对人们所谓的幸福充满憎恨，不得不将那些倾注心血创造出的角色推入悲剧的深渊。”以及稍后那句：“唯有能够唱出人间赞歌的高洁灵魂，才能拯救这个故事。”\cite{ref29}（关于丘比与《萩萩公主》中多罗斯玛亚的可能原型，即 E.T.A. Hoffmann 笔下的朵谢梅（Drosselmeier）之间的更多关联，参见第20章和第21章。）

除其他面向外，《魔法少女小圆》也是一部关于“同意”与“自主权”的动画。到目前为止，剧中已经数次暗示这一主题；而在下一集中，它将变得无可否认。既然如此，还有什么反派比那个控制角色行动者更适合这部作品呢？丘比的所有能力都与此吻合：创造魔法少女，知道她们在想什么，并同时出现在故事各处。它的所有动机——让魔法少女经历情绪的高峰与低谷，让故事世界尽可能长久地运转下去——也同它在不自知中扮演的角色完全一致：丘比就是《魔法少女小圆》的作者。
